\subsection{Orbit dynamics and a force--scattering--delay bridge (interface)}
\label{app:orbit_dynamics_and_force_scattering_bridge}

This appendix extends the unified orbit--gauge--force packaging (Appendix~\ref{app:unified_orbit_gauge_force}) by recording a minimal \emph{particle-orbit dynamics} interface and a minimal \emph{force--phase--delay} bridge.
The statements here are interface-level: they clarify how the paper's CAP-closed action representative and delay dictionaries connect to standard dynamical and scattering observables under explicitly stated prerequisites.

\subsubsection{Orbit dynamics: from a CAP-closed representative action to a particle equation (interface)}
\paragraph{Orbit variables.}
\InterfaceTag An orbit is a base path together with an internal state, $\mathcal{O}=(x_t,\psi_t)$ (Appendix~\ref{app:unified_orbit_gauge_force}).
At the continuum-representative layer, we write a smooth path $x^\mu(\lambda)$ and (when needed) a worldline parameter $\lambda$ (proper time or any convenient affine parameter), together with a covariant transport rule for $\psi$ fixed by a connection.

\paragraph{Force is response; dynamics is variational.}
\InterfaceTag The paper already fixes the response meaning of force (Section~\ref{subsec:force_as_response}) and closes a finite-family representative action $S_{\mathrm{eff}}$ (Section~\ref{app:cap_continuum_action_closure}), with corresponding Euler--Lagrange/field equations (Section~\ref{app:variational_field_equations}).
The minimal additional interface step is the standard point-particle reduction: treat a localized excitation as an effective worldline degree of freedom whose motion follows from varying a reduced action functional along the orbit.

\begin{proposition}[Minimal coupling yields a Lorentz-force orbit equation (standard interface template)]
\label{prop:lorentz_force_from_minimal_coupling_interface}
\InterfaceTag Fix a continuum representative with metric $g_{\mu\nu}$ and a gauge connection one-form $A=A_\mu \dd x^\mu$ with curvature $F=\dd A$ (abelian) or $F=\dd A + A\wedge A$ (non-abelian).
Consider an effective worldline Lagrangian of minimal-coupling form
\[
L(\lambda)=\frac{m}{2}\,g_{\mu\nu}(x)\,\dot x^\mu(\lambda)\dot x^\nu(\lambda)\ +\ q\,A_\mu(x)\,\dot x^\mu(\lambda),
\]
where $\dot x^\mu=\dd x^\mu/\dd\lambda$, $m>0$ is a mass parameter and $q$ a coupling/charge parameter (in the relevant representation).
Then the Euler--Lagrange equations imply an orbit-deflection equation of Lorentz-force type:
\[
m\,\frac{D\dot x^\mu}{\dd\lambda}
\ =\
q\,{F^\mu}_{\ \nu}(x)\,\dot x^\nu,
\]
where $D/\dd\lambda$ is the covariant derivative along the path induced by the Levi--Civita connection of $g$ and the relevant gauge representation.
\end{proposition}

\begin{proof}
This is a standard variational calculation: the kinetic term yields the geodesic operator $D\dot x/\dd\lambda$, and varying the minimal-coupling term yields the curvature contraction $F(\dot x,\cdot)$.
We record it as an interface template that fixes the dictionary: \emph{given} a CAP-selected action representative containing a gauge sector and its minimal couplings, the corresponding orbit equation contains curvature-driven deflection terms.
\end{proof}

\paragraph{Relation to the paper's notion of ``no-force''.}
\InterfaceTag In this dictionary, ``no force'' on the \emph{base orbit} reduces to the geodesic equation (the $q=0$ case or a vanishing curvature sector), while ``no force'' on the \emph{internal state} remains the covariant transport/parallel transport condition $D_t\psi=0$ (Appendix~\ref{app:unified_orbit_gauge_force}).

\subsubsection{Force--phase--delay: a minimal bridge to Wigner--Smith observables (interface)}
\paragraph{Phase--delay dictionary already fixed.}
\InterfaceTag Appendix~\ref{app:time_mass_delay} fixes the unified phase--delay dictionary:
for a one-channel scattering response $S(\omega)=\e^{\iu\delta(\omega)}$ one has
\[
\tau(\omega)=\frac{\dd\delta}{\dd\omega},
\]
and in the multi-channel unitary case the Wigner--Smith operator and its trace are
$Q(\omega)=-\iu S^\dagger \dd S/\dd\omega$ and $\tau_{\mathrm{WS}}(\omega)=\Tr Q(\omega)$ (equations \eqref{eq:wigner_smith_omega}--\eqref{eq:tau_ws_trace_omega}).

\paragraph{A minimal action-to-phase bridge (dictionary).}
\InterfaceTag In semiclassical regimes and/or in effective one-channel reductions, a measured scattering phase can be modeled as a reduced on-shell action (or phase-advance functional) in units of the reference phase scale:
\[
\delta(\omega)\ \approx\ \frac{1}{\hbar}\,S_{\mathrm{red,on\text{-}shell}}(\omega)\ +\ \delta_0,
\]
where $\delta_0$ is a convention-dependent constant.
This is a dictionary statement: it is \emph{not} claimed to be derived from the tick-only folding core.

\begin{proposition}[Force perturbations induce delay perturbations via phase derivatives (interface consequence)]
\label{prop:force_to_delay_via_phase_derivative_interface}
\InterfaceTag Under the phase--action dictionary above, a small perturbation of the effective dynamics that changes the reduced on-shell action by $\Delta S_{\mathrm{red}}(\omega)$ induces a phase change
\[
\Delta\delta(\omega)\approx \frac{1}{\hbar}\,\Delta S_{\mathrm{red}}(\omega),
\]
and therefore a delay change
\[
\Delta\tau(\omega)\approx \frac{1}{\hbar}\,\frac{\dd}{\dd\omega}\Delta S_{\mathrm{red}}(\omega).
\]
Equivalently, in a multi-channel unitary setting, changes in the total scattering phase $\Theta(\omega)=\arg\det S(\omega)$ control changes in $\tau_{\mathrm{WS}}(\omega)=\dd\Theta/\dd\omega$ (Appendix~\ref{app:time_mass_delay}).
\end{proposition}

\paragraph{Scope note.}
\AuditTag This bridge only fixes a minimal closure loop between three already-defined interface objects: force as response (Section~\ref{subsec:force_as_response}), scattering phase/delay (Appendix~\ref{app:time_mass_delay}), and the CAP-selected action representative (Section~\ref{app:cap_continuum_action_closure}).
It does not claim nonperturbative control of interacting 4D scattering, which remains explicitly scoped (Appendix~\ref{app:scattering_haag_ruelle_lsz_interface}).
